\subsection{Observables and Measurements}

\subsubsection{Observables: What Are We Measuring?}

Let's start from the basics: \textbf{\emph{what is an observable?}} In quantum mechanics, an \definition{Observable} is \hl{a mathematical object representing a measurable physical quantity}.
\begin{itemize}
    \item For example, in classical physics, observables are physical quantities that we can measure, such as position, momentum, energy, etc. (e.g., the position of a particle, the energy of a system).
    \item In quantum mechanics, the state is no longer a point, but rather a vector (wavefunction) in a complex vector space (Hilbert space). Because of this, observables can no longer be ordinary functions of the state, but instead must be represented as \textbf{operators} that act on the state vector.
\end{itemize}

\begin{definitionbox}[: Observable]
    In quantum mechanics, an \definition{Observable} is a \textbf{Hermitian operator} $O$ associated with a measurable quantity of a quantum system:
    \begin{equation*}
        O = O^{\dagger}
    \end{equation*}
    Where $O^{\dagger}$ denotes the Hermitian adjoint (conjugate transpose) of $O$.

    \highspace
    With \emph{conjugate transpose}, we mean that if $O$ is represented as a matrix, then $O^{\dagger}$ is obtained by taking the \textbf{transpose} of the matrix and then taking the \textbf{complex conjugate of each element}. The condition $O = O^{\dagger}$ ensures that the eigenvalues of $O$ are real, which is necessary for them to correspond to measurable quantities in quantum mechanics.
\end{definitionbox}

\begin{flushleft}
    \textcolor{Green3}{\faIcon{question-circle} \textbf{Why does quantum mechanics specifically choose Hermitian operators for observables?}}
\end{flushleft}
The answer is simple: \hl{because measurements must produce physically meaningful outcomes.} Specifically, Hermitian operators guarantee: real eigenvalues and orthonormal eigenvectors.
\begin{enumerate}
    \item \important{Real eigenvalues}. The eigenvalues of a Hermitian operator are always real numbers, which is \textbf{essential because measurement outcomes must be real values}. We cannot measure $3 + 2i$ as a physical quantity, but we can measure $3$. Mathematically, if $O$ is a Hermitian operator, then:
    \begin{equation*}
        O\ket{k} = \lambda_k \ket{k}
    \end{equation*}
    Where $\lambda_k$ is the possible measurement outcome (eigenvalue) corresponding to the eigenvector $\ket{k}$, and $\ket{k}$ is the state of the system after the measurement.
    \begin{remarkbox}
        When we write:
        \begin{equation*}
            O\ket{k}=\lambda_k\ket{k}
        \end{equation*}
        We are choosing $\ket{k}$ to be one of the eigenvectors of $O$, and $\lambda_k$ is its corresponding eigenvalue. It is the definition of the eigenvector/eigenvalue relationship.

        \highspace
        Suppose $O$ is an operator, so a matrix. An eigenvector \ket{k} of $O$ is a special vector such that, when $O$ acts on it, the result is the \textbf{same vector direction}, only multiplied by a number. That number is called the eigenvalue $\lambda_k$. So:
        \begin{equation*}
            O\ket{k} = \lambda_k \ket{k}
        \end{equation*}
        Means applying the operator $O$ to \ket{k} does not change its direction; it only scales it by $\lambda_k$.

        \begin{examplebox}[: Eigenvectors and Eigenvalues of the Pauli-Z Operator]
            For example, let's take the Pauli-Z operator:
            \begin{equation*}
                Z =
                \begin{bmatrix}
                    1 & 0 \\
                    0 & -1
                \end{bmatrix}
            \end{equation*}
            And the state \ket{0}:
            \begin{equation*}
                \ket{0} =
                \begin{bmatrix}
                    1 \\
                    0
                \end{bmatrix}
            \end{equation*}
            Applying $Z$ to \ket{0}:
            \begin{equation*}
                Z\ket{0} =
                \begin{bmatrix}
                    1 & 0 \\
                    0 & -1
                \end{bmatrix}
                \begin{bmatrix}
                    1 \\
                    0
                \end{bmatrix}
                =
                \begin{bmatrix}
                    1 \\
                    0
                \end{bmatrix}
                =
                1 \cdot \ket{0}
                =
                \lambda_0 \ket{0}
            \end{equation*}
            Therefore, \ket{0} is an eigenvector of $Z$ with eigenvalue $\lambda_0 = 1$. Similarly, applying $Z$ to \ket{1}:
            \begin{equation*}
                Z\ket{1} =
                \begin{bmatrix}
                    1 & 0 \\
                    0 & -1
                \end{bmatrix}
                \begin{bmatrix}
                    0 \\
                    1
                \end{bmatrix}
                =
                \begin{bmatrix}
                    0 \\
                    -1
                \end{bmatrix}
                =
                -1 \cdot \ket{1}
                =
                \lambda_1 \ket{1}
            \end{equation*}
            Therefore, \ket{1} is an eigenvector of $Z$ with eigenvalue $\lambda_1 = -1$.
        \end{examplebox}

        Beyond the mathematical definition, \ket{0} is a special state of the observable $Z$, measuring $Z$ when the system is in \ket{0} always gives the outcome $+1$. So, when we say that \ket{0} and \ket{1} are eigenvectors of $Z$, we are saying that these are the states in which the measurement of $Z$ is \textbf{perfectly predictable}, yielding the eigenvalues $+1$ and $-1$, respectively.
        
        \newpage

        This is a fundamental aspect, because:
        \begin{itemize}
            \item[\textcolor{Green3}{\faIcon{check}}] If the \textbf{system is already in an eigenvector}: there is no uncertainty, and the measurement outcome is known with certainty.
            \item[\textcolor{Red3}{\faIcon{balance-scale}}] In contrast, if the \textbf{system is not in an eigenvector}: it is a combination of eigenvectors, and measurement becomes \hl{probabilistic}. For example, suppose the qubit is in:
            \begin{equation*}
                \ket{+} = \frac{1}{\sqrt{2}}\left(\ket{0} + \ket{1}\right)
            \end{equation*}
            This is \textbf{not} an eigenvector of $Z$, because applying $Z$:
            \begin{align*}
                Z\ket{+} =&
                \frac{1}{\sqrt{2}}\left(
                    \begin{bmatrix}
                        1 & 0 \\
                        0 & -1
                    \end{bmatrix}
                    \begin{bmatrix}
                        1 \\
                        0
                    \end{bmatrix}
                    +
                    \begin{bmatrix}
                        1 & 0 \\
                        0 & -1
                    \end{bmatrix}
                    \begin{bmatrix}
                        0 \\
                        1
                    \end{bmatrix}
                \right) \\
                =&
                \frac{1}{\sqrt{2}}\left(
                    \begin{bmatrix}
                        1 \\
                        0
                    \end{bmatrix}
                    +
                    \begin{bmatrix}
                        0 \\
                        -1
                    \end{bmatrix}
                \right) \\
                =&
                \frac{1}{\sqrt{2}}\left(
                    \ket{0} - \ket{1}
                \right) \\
                =&
                \ket{-}
            \end{align*}
            \textcolor{Green3}{\faIcon{question-circle} \textbf{Why is \ket{+} not an eigenvector of $Z$?}} If \ket{+} were an eigenvector, we would have by definition:
            \begin{equation*}
                Z\ket{+} = \lambda \ket{+}
            \end{equation*}
            For some eigenvalue $\lambda$. However, as we calculated, applying $Z$ to \ket{+} gives us \ket{-}, which is a different state. Therefore, there is no scalar $\lambda$ such that $Z\ket{+} = \lambda \ket{+}$, confirming that \ket{+} is not an eigenvector of $Z$. Physically, this means that if the system is in state \ket{+}, measuring $Z$ does not yield a definite outcome; instead, it yields $+1$ with probability 0.5 and $-1$ with probability 0.5, reflecting the superposition of eigenstates.
        \end{itemize}
        In summary, if applying an observable $O$ to a state \ket{\psi} produces the same state multiplied by a scalar:
        \begin{equation*}
            O\ket{\psi} = \lambda \ket{\psi}
        \end{equation*}
        Then \ket{\psi} is an eigenvector of $O$, and $\lambda$ is its eigenvalue. In this case, the state already has a definite value for that observable, so measuring $O$ will return $\lambda$ with certainty.

        \highspace
        \textcolor{Green3}{\faIcon{check} \textbf{Physical consequence.}} If the system is in an eigenvector: the \hl{measurement is deterministic}, the \hl{result is known in advance}, the \hl{outcome is the corresponding eigenvalue}.
    \end{remarkbox}
    \begin{proof}[\textbf{Why Hermitian operators have real eigenvalues}]
        Let $A$ be a Hermitian operator:
        \begin{equation*}
            A = A^\dagger
        \end{equation*}
        Assume that \ket{x} is an eigenvector of $A$ with eigenvalue $\lambda$:
        \begin{equation*}
            A\ket{x} = \lambda \ket{x}
        \end{equation*}
        Consider the quantity:
        \begin{equation*}
            \bra{x}A\ket{x}
        \end{equation*}
        Using the eigenvalue equation (i.e., $A\ket{x} = \lambda \ket{x}$) on the right:
        \begin{equation*}
            \bra{x}\underbrace{A\ket{x}}_{\lambda \ket{x}}
            =
            \bra{x}\left(\lambda \ket{x}\right)
            =
            \lambda \braket{x}{x}
        \end{equation*}
        This is possible because $\lambda$ is just a scalar, so it can be moved outside bra-ket notation.\footnote{
            Remember that bras and kets can be represented using vectors and matrices, so the usual rules of linear algebra apply. In particular, since $\lambda$ is a scalar, it can be factored out of an inner product:
            \begin{equation*}
                \bra{x}(\lambda \ket{x})
                =
                \lambda \braket{x}{x}
            \end{equation*}
            More generally, scalar multiplication commutes with matrix products. Therefore, for matrices $A$ and $B$:
            \begin{equation*}
                \lambda AB = A\lambda B = AB\lambda
            \end{equation*}
            Provided that the order of the matrices themselves is preserved.
        }
        Now take the Hermitian conjugate of the eigenvalue equation:
        \begin{equation*}
            A\ket{x} = \lambda \ket{x}
            \implies
            \bra{x}A^\dagger = \overline{\lambda} \bra{x}
        \end{equation*}
        Since $A$ is Hermitian:
        \begin{equation*}
            A^\dagger = A
        \end{equation*}
        Therefore:
        \begin{equation*}
            \bra{x}A = \overline{\lambda} \bra{x}
        \end{equation*}
        Applying this on the left of $\bra{x} A\ket{x}$:
        \begin{equation*}
            \underbrace{\bra{x}A}_{\overline{\lambda} \bra{x}}\ket{x}
            =
            \left(\overline{\lambda}\bra{x}\right)\ket{x}
            =
            \overline{\lambda}\braket{x}{x}
        \end{equation*}
        Combining the two expressions:
        \begin{equation*}
            \lambda \braket{x}{x}
            =
            \overline{\lambda}\braket{x}{x}
        \end{equation*}
        Since:
        \begin{equation*}
            \braket{x}{x} > 0
        \end{equation*}
        Because the inner product of a vector with itself is always positive (unless the vector is the zero vector, which cannot be an eigenvector), we can divide both sides by $\braket{x}{x}$:
        \begin{equation*}
            \lambda = \overline{\lambda}
        \end{equation*}
        Therefore, $\lambda$ is real.
    \end{proof}

    
    \newpage


    \item \important{Orthonormal eigenvectors}. A vector is:
    \begin{itemize}
        \item \textbf{Orthogonal} if its inner product with another vector is zero (i.e., they are perpendicular in the vector space). For example, in a 2D space, the vectors $\ket{0}$ and $\ket{1}$ are orthogonal because:
        \begin{equation*}
            \braket{0}{1}
            =
            \begin{bmatrix}
                1 & 0
            \end{bmatrix}
            \begin{bmatrix}
                0 \\
                1
            \end{bmatrix}
            = 0
        \end{equation*}

        \item \textbf{Normalized} if its inner product with itself is equal to one (i.e., it has unit length). For example:
        \begin{equation*}
            \braket{0}{0}
            =
            \begin{bmatrix}
                1 & 0
            \end{bmatrix}
            \begin{bmatrix}
                1 \\
                0
            \end{bmatrix}
            = 1
        \end{equation*}
    \end{itemize}

    This property is \hl{crucial because quantum measurement requires a basis whose states are both perfectly distinguishable (orthogonal) and probabilistically consistent (normalized).}
    \begin{itemize}
        \item[\textcolor{Red2}{\faIcon{times}}] \textbf{If eigenvectors were not orthogonal}, different measurement\break states could partially overlap:
        \begin{equation*}
            \braket{v_1}{v_2} \neq 0
        \end{equation*}
        In this case, the measurement process would not be able to distinguish the states perfectly, leading to ambiguity in the interpretation of measurement outcomes.
        \begin{examplebox}[: Why Orthogonality Matters]
            Consider the two states:
            \begin{equation*}
                \ket{v_1} = \ket{0}
                \qquad
                \ket{v_2} = \frac{1}{\sqrt{2}}(\ket{0} + \ket{1})
            \end{equation*}

            Their inner product is:
            \begin{equation*}
                \braket{v_1}{v_2}
                =
                \bra{0}
                \frac{1}{\sqrt{2}}(\ket{0} + \ket{1})
                =
                \frac{1}{\sqrt{2}}
            \end{equation*}

            Since:
            \begin{equation*}
                \braket{v_1}{v_2} \neq 0
            \end{equation*}
            the two states are not orthogonal. This means that the \textbf{states partially overlap}, and therefore a \textbf{measurement cannot distinguish them perfectly}. For example, if the system is in state $\ket{v_2}$, there is still a non-zero probability of confusing it with $\ket{v_1}$ during a measurement. Therefore, the measurement outcomes would be ambiguous, and we would not be able to assign a clear physical meaning to the measurement results.
        \end{examplebox}

        \item[\textcolor{Red2}{\faIcon{times}}] \textbf{If eigenvectors were not normalized}:
        \begin{equation*}
            \braket{v}{v} \neq 1
        \end{equation*}
        Then the Born rule (\autopageref{eq:born-rule}) would not produce properly normalized probabilities, making the probabilistic interpretation of quantum mechanics inconsistent.
        \begin{examplebox}[: Why Normalization Matters]
            Consider the state:
            \begin{equation*}
                \ket{v}
                =
                \begin{pmatrix}
                    2 \\
                    0
                \end{pmatrix}
            \end{equation*}

            Its norm is:
            \begin{equation*}
                \braket{v}{v}
                =
                \begin{bmatrix}
                    2 & 0
                \end{bmatrix}
                \begin{bmatrix}
                    2 \\
                    0
                \end{bmatrix}
                = 4
            \end{equation*}

            Therefore, the state is not normalized.

            \highspace
            If we tried to interpret the amplitudes probabilistically using the Born rule, we would obtain:
            \begin{equation*}
                p_0 = |2|^2 = 4
            \end{equation*}
            Which is \textbf{impossible}, since \textbf{probabilities must lie between 0 and 1}, and their total sum must equal 1.

            \highspace
            To obtain a valid quantum state, we must normalize the vector:
            \begin{equation*}
                \ket{v_{\text{norm}}}
                =
                \frac{1}{2}
                \begin{pmatrix}
                    2 \\
                    0
                \end{pmatrix}
                =
                \begin{pmatrix}
                    1 \\
                    0
                \end{pmatrix}
            \end{equation*}
            Therefore, without normalization, we would end up with an invalid quantum state that cannot be interpreted probabilistically, and the measurement outcomes would not correspond to any physical reality.
        \end{examplebox} 
    \end{itemize}
    \begin{proof}[\textbf{Why eigenvectors of Hermitian operators are orthogonal}]
        Let $A$ be Hermitian and let:
        \begin{equation*}
            A\ket{x} = \lambda \ket{x}, \qquad A\ket{y} = \mu \ket{y}
        \end{equation*}
        Where $\lambda \ne \mu$ are the eigenvalues corresponding to the eigenvectors $\ket{x}$ and $\ket{y}$, respectively. Consider the inner product:
        \begin{equation*}
            \bra{y} A \ket{x}
        \end{equation*}
        Because we want to show that $\ket{x}$ and $\ket{y}$ are orthogonal, we need to show that $\braket{y}{x} = 0$.
        Using the eigenvalue equation on the right:
        \begin{equation*}
            \bra{y} A \ket{x}
            =
            \bra{y} (\lambda \ket{x})
            =
            \lambda \braket{y}{x}
        \end{equation*}
        We moved $\lambda$ outside the inner product because it is a scalar (see the footnote in the previous proof for more details on this step).
        Now use the Hermitian property on the left:
        \begin{equation*}
            \underbrace{\bra{y} A }_{\bra{y} A^\dagger} = \mu \bra{y}
        \end{equation*}
        Therefore:
        \begin{equation*}
            \bra{y} A \ket{x}
            =
            \mu \braket{y}{x}
        \end{equation*}
        Combining the two expressions:
        \begin{equation*}
            \lambda \braket{y}{x}
            =
            \mu \braket{y}{x}
        \end{equation*}
        Thus:
        \begin{equation*}
            \left(\lambda - \mu\right) \braket{y}{x} = 0
        \end{equation*}
        Since $\lambda \neq \mu$ by assumption, we must have:
        \begin{equation*}
            \braket{y}{x} = 0
        \end{equation*}
        Therefore, the eigenvectors are orthogonal.
    \end{proof}
\end{enumerate}
We can think of a \hl{Hermitian operator} as a sort of ``\emph{machine}'' that \hl{decomposes the quantum state into measurable components} (eigenvectors) \hl{and assigns a real value} (eigenvalue) \hl{to each component}, which corresponds to the possible outcomes of a measurement. The orthonormality of the eigenvectors ensures that these outcomes are distinct and can be interpreted probabilistically in a consistent way. In summary, a Hermitian observable tells us:
\begin{itemize}
    \item \emph{What outcomes are possible when we measure a quantum system.}
    \item \emph{With respect to which basis we are measuring the system.}
    \item \emph{Into which states the system will collapse after the measurement.}
\end{itemize}

\highspace
\begin{flushleft}
    \textcolor{Green3}{\faIcon{book} \textbf{Spectral Theorem}}
\end{flushleft}
So far, we know that an observable $O$ is a Hermitian operator, it has eigenvectors \ket{k}, and each eigenvector has an eigenvalue $\lambda_k$. So if the system is in \ket{k}, measuring $O$ returns $\lambda_k$ with certainty. Now, a natural question arises: \textbf{what happens if the system is in an arbitrary state \ket{k}, not necessarily one of the eigenvectors?} To answer this question, the \hl{Spectral Theorem} comes into play by showing that \hl{every observable can be written as a sum of its eigenvectors and eigenvalues}.

\highspace
Suppose the observable has a set of eigenvectors \ket{k} with corresponding eigenvalues $\lambda_k$. Since the eigenvectors form an orthonormal basis, any quantum state \ket{v} can be expressed as a linear combination of these eigenvectors (\autopageref{eq:general-multiple-qubit-state}):
\begin{equation*}
    \ket{v} = \sum_k c_k \ket{k}
\end{equation*}
Where $c_k$ are complex coefficients. The \definition{Spectral Theorem} states that the observable $O$ can be expressed as:
\begin{equation}\label{eq:spectral-decomposition}
    O = \sum_k \lambda_k \proj{k}
\end{equation}
Where each term $\lambda_k \proj{k}$ contains two pieces of information:
\begin{itemize}
    \item \important{Projector \proj{k}} (\autopageref{eq:measurement-operator}) onto the eigenvector \ket{k}, which extracts the component of the state along \ket{k} during measurement. It is defined as:
    \begin{equation*}
        \proj{k} = \ket{k}\bra{k}
    \end{equation*}
    \item \important{Eigenvalue $\lambda_k$}, which assigns the measurement outcome associated with that component.
\end{itemize}
\textcolor{Green3}{\faIcon{question-circle} \textbf{Why the Spectral Theorem is important?}} Suppose someone gives us an observable $O$ without telling us its eigenvectors or eigenvalues:
\begin{equation*}
    O =
    \begin{bmatrix}
        2 & 1 \\
        1 & 2
    \end{bmatrix}
\end{equation*}
Without the Spectral Theorem, the physical meaning of the observable would not be immediately apparent. However, by applying the Spectral Theorem, we can decompose $O$ into its eigenvectors and eigenvalues, allowing us to \hl{understand the measurement process and predict the outcomes based on the state of the system}. For example, we can find the eigenvalues and eigenvectors of $O$:
\begin{itemize}
    \item To find the eigenvalues, we solve the characteristic equation:
    \begin{align*}
        \det\left(O - \lambda I\right) &= 0 \\
        \det\left(
            \begin{bmatrix}
                2 - \lambda & 1 \\
                1 & 2 - \lambda
            \end{bmatrix}
        \right) &= 0 \\
        \left(2 - \lambda\right)\left(2 - \lambda\right) - 1 \cdot 1 &= 0 \\
        4 - 4\lambda + \lambda^2 - 1 &= 0 \\
        \lambda^2 - 4\lambda + 3 &= 0 \\
        \left(\lambda - 3\right)\left(\lambda - 1\right) &= 0
    \end{align*}
    Thus, the eigenvalues are $\lambda_1 = 3$ and $\lambda_2 = 1$.
    \item Eigenvectors can be found by solving $(O - \lambda I)\ket{k} = 0$ for each eigenvalue:
    \begin{itemize}
        \item For $\lambda_1 = 3$:
        \begin{align*}
            \left(O - 3I\right)\ket{v_1} &= 0 \\
            \left(\begin{bmatrix}
                2 - 3 & 1 \\
                1 & 2 - 3
            \end{bmatrix}\right)\ket{v_1} &= 0 \\
            \left(\begin{bmatrix}
                -1 & 1 \\
                1 & -1
            \end{bmatrix}\right)\ket{v_1} &= 0
        \end{align*}
        Now to find the eigenvector, we need to solve the system of equations:
        \begin{equation*}
            \begin{bmatrix}
                -1 & 1 \\
                1 & -1
            \end{bmatrix} \begin{bmatrix}
                x \\
                y
            \end{bmatrix}
            =
            \begin{bmatrix}
                0 \\
                0
            \end{bmatrix}
        \end{equation*}
        That is a homogeneous system of linear equations:
        \begin{equation*}
            \begin{cases}
                -x + y = 0 \\
                x - y = 0
            \end{cases}
            =
            \begin{cases}
                y = x \\
                x = y
            \end{cases}
        \end{equation*}
        So any vector of the form $\left(x, x\right)$ is an eigenvector. We choose the simplest one $\left(1, 1\right)$, and then we normalize it:
        \begin{equation*}
            \left\| \begin{bmatrix}
                1 \\
                1
            \end{bmatrix} \right\| = \sqrt{1^2 + 1^2} = \sqrt{2}
        \end{equation*}
        Therefore, the normalized eigenvector is:
        \begin{equation*}
            \ket{v_1} = \frac{1}{\sqrt{2}}\left(\ket{0} + \ket{1}\right)
        \end{equation*}

        \item For $\lambda_2 = 1$:
        \begin{align*}
            \left(O - 1I\right)\ket{v_2} &= 0 \\
            \left(\begin{bmatrix}
                2 - 1 & 1 \\
                1 & 2 - 1
            \end{bmatrix}\right)\ket{v_2} &= 0 \\
            \left(\begin{bmatrix}
                1 & 1 \\
                1 & 1
            \end{bmatrix}\right)\ket{v_2} &= 0
        \end{align*}
        We need to solve the system of equations:
        \begin{equation*}
            \begin{cases}
                x + y = 0 \\
                x + y = 0
            \end{cases}
            \Rightarrow
            x + y = 0 \Rightarrow y = -x \xrightarrow{\text{substitute from }\begin{bmatrix}
                x \\ y
            \end{bmatrix}}
            \begin{bmatrix}
                x \\
                -x
            \end{bmatrix}
        \end{equation*}
        So any vector of the form $\left(x, -x\right)$ is an eigenvector. We choose the simplest one $\left(1, -1\right)$, and then we normalize it:
        \begin{equation*}
            \left\| \begin{bmatrix}
                1 \\
                -1
            \end{bmatrix} \right\| = \sqrt{1^2 + (-1)^2} = \sqrt{2}
        \end{equation*}
        Therefore, the normalized eigenvector is:
        \begin{equation*}
            \ket{v_2} = \frac{1}{\sqrt{2}}\left(\ket{0} - \ket{1}\right)
        \end{equation*}
    \end{itemize}
\end{itemize}
Then we can express $O$ in terms of its spectral decomposition:
\begin{equation*}
    O = 3 \proj{v_1} + 1 \proj{v_2}
\end{equation*}
\hl{From this decomposition, we can immediately read off the possible outcomes} from the eigenvalues ($3$ and $1$), \hl{the measurement basis} from the eigenvectors ($\ket{v_1}$ and $\ket{v_2}$), and \hl{the measurement operators} from the projectors ($\proj{v_1}$ and $\proj{v_2}$). It also allows us to calculate the probabilities of obtaining each outcome when measuring $O$ on an arbitrary state \ket{\psi} using the Born rule:
\begin{equation*}
    p_1 = |\braket{v_1}{\psi}|^2
    \qquad
    p_2 = |\braket{v_2}{\psi}|^2
\end{equation*}
Where $p_1$ and $p_2$ are the probabilities of obtaining the outcomes $3$ and $1$, respectively, when measuring the observable $O$ on a state \ket{\psi}. And finally the expectation value of $O$ in the state \ket{\psi} can be calculated as:
\begin{equation*}
    \bra{\psi} O \ket{\psi} = \sum_k \lambda_k p_k
\end{equation*}
In this example, this becomes:
\begin{equation*}
    \bra{\psi} O \ket{\psi} = \sum_{k=1}^2 \lambda_k p_k = 3p_1 + 1p_2
\end{equation*}
In summary, the \hl{Spectral Theorem transforms an abstract Hermitian matrix into a complete description of a quantum measurement.} It reveals:
\begin{itemize}
    \item The \textbf{possible measurement outcomes} (eigenvalues),
    \item The \textbf{corresponding post-measurement states} (eigenvectors),
    \item The \textbf{projectors used to compute probabilities},
    \item And the \textbf{expectation value} as a weighted average of the outcomes.
\end{itemize}